O \textit{matching} já foi amplamente utilizado em diferentes áreas de políticas públicas como uma estratégia central de avaliação causal em contextos nos quais experimentos aleatorizados não são viáveis. Em muitas políticas reais, a randomização esbarra em restrições éticas (não é aceitável negar um serviço potencialmente benéfico), políticas (programas são definidos por critérios legais ou administrativos) ou operacionais (implementação gradual, descentralizada ou baseada em demanda). Nesses ambientes, avaliações baseadas exclusivamente em comparações ingênuas entre participantes e não participantes tendem a produzir estimativas fortemente enviesadas.

O \textit{matching} surge exatamente como uma resposta a esse problema. Ele operacionaliza a intuição de ``comparar semelhantes com semelhantes'', mas o faz de forma sistemática e transparente, ancorada em um arcabouço formal de inferência causal. Ao explicitar as covariadas utilizadas, o critério de comparabilidade e as regiões do espaço de dados nas quais a inferência é possível, o método busca aproximar o desenho observacional das condições de um experimento controlado, deixando claras tanto suas virtudes quanto suas limitações \citep{rosenbaum1983central, rubin1974estimating}.

\subsection*{Educação}

Na área de educação, o \textit{matching} tem sido amplamente empregado na avaliação de programas de educação complementar, como cursos técnicos, programas de qualificação profissional, bolsas de estudo, tutorias acadêmicas e iniciativas de reforço escolar. Esses programas raramente são alocados aleatoriamente: a participação costuma depender de características individuais, critérios administrativos, decisões familiares ou da própria escola, o que gera forte seleção dos participantes.

Avaliações baseadas em \textit{matching} constroem grupos de comparação a partir de indivíduos não participantes com perfil semelhante em termos de idade, escolaridade prévia, desempenho acadêmico anterior, renda familiar, características da escola e localização geográfica. Essa abordagem permite estimar efeitos causais sobre uma ampla gama de desfechos, incluindo progressão educacional, evasão escolar, desempenho em testes padronizados, inserção no mercado de trabalho e salários iniciais \citep{heckman1997matching, smith2005does}.

Um conjunto particularmente influente de estudos comparou estimativas obtidas via \textit{matching} com resultados de experimentos aleatorizados em programas de treinamento educacional e profissional nos Estados Unidos. Esses trabalhos mostraram que, quando o conjunto de covariadas é rico e o suporte comum é adequadamente respeitado, o \textit{matching} é capaz de reproduzir de forma bastante próxima os efeitos estimados experimentalmente, fornecendo evidência empírica de sua utilidade prática em contextos educacionais \citep{dehejia1999causal, dehejia2002propensity}.

\subsection*{Saúde}

Na área da saúde, o \textit{matching} é frequentemente utilizado para avaliar programas de prevenção, acompanhamento médico e intervenções em saúde pública, especialmente quando experimentos são inviáveis ou controversos. Ao analisar, por exemplo, o impacto de campanhas de vacinação, programas de rastreamento precoce ou acompanhamento preventivo, os participantes tendem a diferir sistematicamente dos não participantes não apenas em características observáveis — como idade, sexo e histórico clínico —, mas também em comportamentos relacionados ao autocuidado, adesão a tratamentos e aversão ao risco.

Esse padrão de autoseleção gera viés de seleção substancial. Indivíduos mais atentos à própria saúde tendem simultaneamente a participar das intervenções e a apresentar melhores desfechos, mesmo na ausência do programa. O \textit{matching} contribui para mitigar esse viés ao construir grupos comparáveis em termos de histórico médico, utilização prévia de serviços de saúde e perfil demográfico, aproximando a análise do contrafactual que seria observado em um experimento controlado \citep{stuart2010matching}.

Em aplicações baseadas em grandes bases administrativas e registros hospitalares, o \textit{matching} desempenha um papel fundamental ao tornar esses dados observacionais utilizáveis para inferência causal. Ao mesmo tempo, esses estudos reforçam um ponto central do método: ele corrige apenas o viés associado às variáveis observáveis. Fatores não observados — como preferências individuais, informação imperfeita ou atitudes em relação ao risco — permanecem potenciais fontes de viés, exigindo cautela na interpretação dos resultados \citep{imbens2015causal}.

\subsection*{Mercado de Trabalho}

No mercado de trabalho, técnicas de \textit{matching} ocupam um papel central na avaliação de políticas ativas de emprego, como subsídios à contratação, programas de intermediação de mão de obra, seguro-desemprego com condicionalidades e iniciativas de qualificação profissional. Esses programas são tipicamente implementados com critérios administrativos ou regras de elegibilidade que geram forte seleção dos participantes, inviabilizando comparações diretas simples.

Avaliações empíricas utilizam \textit{matching} para comparar beneficiários do programa com trabalhadores não beneficiários de perfil semelhante em termos de idade, escolaridade, ocupação anterior, histórico de emprego, duração do desemprego e região de residência. Essa estratégia permite identificar se a política gerou efeitos sobre a probabilidade de inserção no emprego, a estabilidade ocupacional, a duração do desemprego ou os salários \citep{heckman1998characterizing, caliendo2008some}.

Esses estudos frequentemente exploram dados administrativos longitudinais, que oferecem informação detalhada sobre trajetórias no mercado de trabalho, mas carecem de uma estratégia experimental explícita. O \textit{matching} torna-se, assim, uma ferramenta-chave para transformar esses dados em evidência causal plausível, explicitando as hipóteses sob as quais as comparações são válidas e revelando as regiões do espaço de dados em que a inferência é possível \citep{lechner2011matching}.

\subsection*{Boas práticas e limitações}

A ampla aplicação do \textit{matching} em políticas públicas também levou ao desenvolvimento de um conjunto de boas práticas empíricas. Entre elas, destacam-se a verificação sistemática de balanceamento das covariadas após o pareamento, a análise explícita do suporte comum, a avaliação da sensibilidade dos resultados a diferentes especificações do método e o uso de procedimentos apropriados de inferência estatística \citep{abadie2006large}.

Ao mesmo tempo, a literatura enfatiza que o \textit{matching} não elimina a necessidade de julgamento substantivo. A escolha das covariadas, a definição da população de interesse e a interpretação dos resultados continuam sendo decisões centrais do pesquisador. O método não cria informação onde ela não existe: quando não há sobreposição suficiente entre tratados e controles, a consequência lógica é a exclusão de observações e a redefinição do parâmetro causal estimado.

\subsection*{Resumo}

Os exemplos discutidos mostram que o \textit{matching} não é apenas uma ferramenta acadêmica, mas uma \textbf{estratégia prática de avaliação}, amplamente utilizada em áreas centrais como educação, saúde e mercado de trabalho. Seu papel fundamental é ampliar a credibilidade das comparações em contextos observacionais, tornando mais plausível a interpretação causal dos resultados quando a experimentação não é possível.

Embora não substitua experimentos aleatorizados, o \textit{matching} expande de forma substantiva o conjunto de políticas públicas que podem ser avaliadas de maneira crível. Ao tornar explícitas as hipóteses identificadoras, as decisões de desenho e as limitações impostas pelos dados, o método contribui para avaliações mais transparentes, replicáveis e informativas em ambientes empíricos complexos.
